\documentclass[11pt]{article}

\usepackage[utf8]{inputenc}
\usepackage[T1]{fontenc}
\usepackage{lmodern}
\usepackage{amsmath,amssymb,amsthm,mathtools}
\usepackage{stmaryrd}
\usepackage[margin=1in]{geometry}
\usepackage{enumitem}
\usepackage{xcolor}
\usepackage{hyperref}
\hypersetup{
  colorlinks=true,
  linkcolor=blue!60!black,
  citecolor=blue!60!black,
  urlcolor=blue!60!black
}

% theorem environments
\newtheorem{theorem}{Theorem}
\newtheorem{lemma}[theorem]{Lemma}
\newtheorem{proposition}[theorem]{Proposition}
\newtheorem{corollary}[theorem]{Corollary}
\newtheorem{question}[theorem]{Question}
\newtheorem{conjecture}[theorem]{Conjecture}
\theoremstyle{definition}
\newtheorem{definition}[theorem]{Definition}
\newtheorem{example}[theorem]{Example}
\newtheorem{remark}[theorem]{Remark}

\usepackage[nameinlink,capitalize]{cleveref}
\crefname{theorem}{Theorem}{Theorems}
\crefname{lemma}{Lemma}{Lemmas}
\crefname{proposition}{Proposition}{Propositions}
\crefname{corollary}{Corollary}{Corollaries}
\crefname{definition}{Definition}{Definitions}
\crefname{remark}{Remark}{Remarks}
\Crefname{theorem}{Theorem}{Theorems}
\Crefname{lemma}{Lemma}{Lemmas}
\Crefname{proposition}{Proposition}{Propositions}
\Crefname{corollary}{Corollary}{Corollaries}
\Crefname{definition}{Definition}{Definitions}
\Crefname{remark}{Remark}{Remarks}

% macros
\newcommand{\R}{\mathbb{R}}
\newcommand{\N}{\mathbb{N}}
\newcommand{\E}{\mathbb{E}}
\newcommand{\Prob}{\mathbb{P}}
\newcommand{\A}{\mathcal{A}}
\newcommand{\F}{\mathcal{F}}
\newcommand{\K}{\mathcal{K}}
\newcommand{\T}{\mathcal{T}}
\newcommand{\Q}{\mathcal{Q}}
\newcommand{\Sset}{\mathcal{S}}
\newcommand{\U}{\mathcal{U}}
\newcommand{\X}{\mathcal{X}}
\newcommand{\C}{\mathcal{C}}
\newcommand{\Homp}{\operatorname{Hom}^{+}}
\newcommand{\Homalg}{\operatorname{Hom}}
\newcommand{\Ext}{\operatorname{Ext}}
\newcommand{\supp}{\operatorname{supp}}
\newcommand{\Forb}{\operatorname{Forb}}
\newcommand{\brak}[1]{\left\langle #1\right\rangle}
\newcommand{\avg}[2]{\left\llbracket #1 \right\rrbracket_{#2}}
\newcommand{\qmap}{\mathfrak q}
\newcommand{\eps}{\varepsilon}
\newcommand{\one}{\mathbf 1}
\DeclareMathOperator{\cl}{cl}

\title{Root Planting, Quotient Semantics, and a Counterexample\\
       for Forbidden-Subgraph Reasoning in Flag Algebras}
\author{}
\date{\today}

\begin{document}
\maketitle

\begin{abstract}
Forbidden-subgraph reasoning in Razborov's flag algebras can be carried out in two
ways: in the \emph{quotient} flag algebra of the constrained theory, or by
\emph{ensemble} semantics inside the ambient theory, asking that an inequality hold
almost surely along the random labelled extensions of constrained unlabelled limits.
When do these two semantics agree?  We show that the answer is not governed merely
by heredity of the constrained class, but by a support-closure condition: every
constrained labelled homomorphism must be approximable by support points of random
labelled extensions of constrained unlabelled homomorphisms.  We call this condition
\emph{root-plantability} and prove that it is exactly equivalent to the
quotient/ensemble correspondence for all flag-algebra elements.  Closure under
independent blow-ups is shown to be a sufficient condition for root-plantability,
which recovers the correspondence for clique-free and triangle-free theories.
Finally, we show that heredity alone is not enough, via a general degeneracy
obstruction: every edge-degenerate hereditary class containing arbitrarily large
stars---for example $C_4$-free, $K_{s,t}$-free, $C_{2k}$-free, and planar
graphs---fails to be root-plantable at the one-vertex type.  The witness is the
rooted edge flag $e$: its negation $-e$ is non-negative under every ensemble
extension yet negative on the centre-rooted star homomorphism in the quotient
theory.  Since root-plantability is preserved by complementation, the same failure
occurs for the dense complement classes; so neither sparsity nor density is
decisive, and the genuine obstruction is that a rooted density is pinned to a
boundary value across all constrained limits.  We isolate this as a general pinning
obstruction and conjecture that its absence characterises root-plantability, with
the $C_5$-free graphs as the critical test case, for which we establish the
supporting local structure.
\end{abstract}

\section{Introduction}

A recurring task in extremal combinatorics is to prove a density inequality that
holds only under a forbidden-subgraph constraint, such as ``in every triangle-free
graph this combination of subgraph densities is non-negative.''  Razborov's flag
algebras offer two routes to such statements.  One may pass to the \emph{quotient}
flag algebra of the constrained theory and ask for non-negativity there
(\emph{quotient semantics}); or one may stay in the ambient theory of all graphs and
ask that the inequality hold almost surely along the random labelled extensions of
constrained limits (\emph{ensemble semantics}).  The first is the algebra one
actually computes in; the second is what a constrained density bound asserts.  When
the two agree, a certificate proved in the quotient is a genuine density bound for
the constrained class.  This paper pins down exactly when they agree.

One direction is automatic.  Quotient semantics always implies ensemble semantics
(Theorem~\ref{thm:support-criterion}), so a quotient certificate is unconditionally
\emph{sound}: the validity of forbidden-subgraph flag-algebra reasoning is never in
question.  The substantive content is the converse---whether \emph{every} true
constrained density inequality can be certified in the quotient---so our results
concern the \emph{completeness} of the quotient calculus, and the counterexamples
below exhibit its \emph{incompleteness} for certain classes.

Let $\T_0$ be a universal first-order theory in a finite relational language,
and let $\T_1$ be a universal strengthening of $\T_0$ in the same language.  In
applications one should think of $\T_0$ as the theory of finite simple graphs and
of $\T_1$ as the theory obtained by imposing a forbidden-subgraph condition.  For
each non-degenerate type $\sigma$ of $\T_1$, Razborov's construction gives a
quotient homomorphism
\[
  \qmap_\sigma:\A^\sigma[\T_0]\twoheadrightarrow \A^\sigma[\T_1].
\]
The quotient semantics of an element $f\in\A^\sigma[\T_0]$ says that
\[
  \psi(\qmap_\sigma f)\ge 0
  \qquad\text{for every }\psi\in\Homp(\A^\sigma[\T_1],\R).
\]
The ensemble semantics stays in the ambient theory $\T_0$.  It asks that, for
every unlabelled $\T_1$-supported homomorphism $\phi_0$, the canonical random
extension $\phi^\sigma$ of $\phi_0$ to type $\sigma$ satisfies
\[
  \phi^\sigma(f)\ge 0
  \qquad\text{almost surely.}
\]

A natural sufficient condition for the equivalence of these two semantics is that
the constrained class be hereditary and closed under independent blow-ups.  Below we
isolate the exact structural reason this works, give a general criterion that is both
necessary and sufficient, and record a simple counterexample showing that heredity
alone does not suffice.

\paragraph{Main message.}
The right property is not simply heredity.  It is the following support
condition:
\[
  \overline{\bigcup_{\phi_0}\supp(\Ext_\sigma(\phi_0))}
  \;=\;
  \{\psi\circ\qmap_\sigma:\psi\in\Homp(\A^\sigma[\T_1],\R)\},
\]
both sides regarded inside the ambient homomorphism space
$\Homp(\A^\sigma[\T_0],\R)$.  Informally, every constrained labelled limit must be
approximable by positive-probability rooted views of constrained unlabelled limits.

\section{Flag-algebra background}

We use Razborov's notation and terminology \cite{Razborov2007}.  For a universal
theory $\T$ and a type $\sigma$, $\A^\sigma[\T]$ denotes the flag algebra of
$\sigma$-flags.  The compact space of positive homomorphisms is
\[
  \Homp(\A^\sigma[\T],\R).
\]
A sequence of finite $\sigma$-flags converges if all flag densities converge; its
limit is an element of this homomorphism space, and conversely every positive
homomorphism arises in this way \cite[Theorem 3.3]{Razborov2007}.  We write
$\F^\sigma[\T]$ for the (countable) set of finite $\sigma$-flags of $\T$ up to
isomorphism, abbreviated $\F^\sigma$ when $\T$ is clear.  Since each flag density
lies in $[0,1]$, evaluating a homomorphism on all flags embeds
$\Homp(\A^\sigma[\T],\R)$ as a closed subspace of the compact metrizable cube
$[0,1]^{\F^\sigma[\T]}$.

For the empty type we write $\A^0[\T]$.  The averaging (unlabelling) operator
\[
  \avg{\cdot}{\sigma}:\A^\sigma[\T]\to\A^0[\T]
\]
forgets the labels of a $\sigma$-flag and averages over the choice of labelling.  The
only properties we use are that it is linear and that on a single flag $F$ with
underlying unlabelled structure $M$ it returns a \emph{non-negative scalar multiple}
of $M$.  Write $1_\sigma\in\A^\sigma[\T]$ for the flag consisting of the type alone
(the multiplicative unit of $\A^\sigma[\T]$).  If $\sigma$ has size $k$, then
\[
  \brak{\sigma}:=\avg{1_\sigma}{\sigma}\in\A^0[\T]
\]
is the unlabelled density of embeddings of the type $\sigma$.  The type $\sigma$ is
\emph{non-degenerate} if $\brak\sigma\neq0$, equivalently if some limit assigns it
positive density.  If
$\phi_0\in\Homp(\A^0[\T],\R)$ and $\phi_0(\brak\sigma)>0$, Razborov's extension
theorem gives a unique probability measure on
$\Homp(\A^\sigma[\T],\R)$, denoted here by
\[
  \Ext_\sigma(\phi_0),
\]
characterised by
\begin{equation}\label{eq:extension-expectation}
  \E_{\phi^\sigma\sim\Ext_\sigma(\phi_0)}[\phi^\sigma(g)]
  =
  \frac{\phi_0(\avg{g}{\sigma})}{\phi_0(\brak\sigma)}
  \qquad(g\in\A^\sigma[\T]).
\end{equation}
Moreover, if a finite sequence converges to $\phi_0$, then the corresponding
finite random-root distributions converge weakly to $\Ext_\sigma(\phi_0)$
\cite[Theorems 3.5 and 3.12]{Razborov2007}.

We will repeatedly use the following elementary consequence of the definition of
support.

\begin{lemma}[Almost-sure non-negativity and support]\label{lem:support-as}
Let $P$ be a Borel probability measure on a compact space $X$, and let
$h:X\to\R$ be continuous.  Then
\[
  P[h\ge0]=1
\]
if and only if $h(x)\ge0$ for every $x\in\supp(P)$.
\end{lemma}

\begin{proof}
If some $x\in\supp(P)$ satisfies $h(x)<0$, then the open set
$\{y:h(y)<0\}$ contains $x$ and hence has positive $P$-measure, a
contradiction.  Conversely, a Borel probability measure is concentrated on its
closed support, so non-negativity on the support implies almost-sure
non-negativity.
\end{proof}

\section{Universal strengthenings and supported homomorphisms}

Let $\T_1$ be a universal strengthening of $\T_0$ in the same relational
language.  For a non-degenerate type $\tau$ of $\T_1$, define
\[
  X_\tau:=\Homp(\A^\tau[\T_0],\R).
\]
The quotient map
\[
  \qmap_\tau:\A^\tau[\T_0]\twoheadrightarrow\A^\tau[\T_1]
\]
pulls every homomorphism on $\A^\tau[\T_1]$ back to a homomorphism on
$\A^\tau[\T_0]$.  We write
\[
  Q_\tau
  :=
  \{\psi\circ\qmap_\tau:
      \psi\in\Homp(\A^\tau[\T_1],\R)\}
  \subseteq X_\tau.
\]
Thus $Q_\tau$ is the constrained, or $\T_1$-supported, part of the ambient
homomorphism space.  It is closed in $X_\tau$: it is cut out by the conditions
that all forbidden finite $\tau$-flags have value zero.

Equivalently, $\chi\in X_\tau$ lies in $Q_\tau$ iff it vanishes on every
$\tau$-flag which is a model of $\T_0$ but not a model of $\T_1$.  This
formulation uses the fact that $\T_1$ is universal: if a finite flag violates
$\T_1$, then its underlying finite structure is one of the forbidden finite
models of the quotient theory.

\begin{lemma}[Support passes to random extensions]\label{lem:support-passes-general}
Let $\phi_0\in Q_0$ and suppose $\phi_0(\brak\sigma)>0$.  If
$\phi^\sigma\sim\Ext_\sigma(\phi_0)$ is formed in the ambient theory $\T_0$,
then
\[
  \phi^\sigma\in Q_\sigma
  \qquad\text{almost surely.}
\]
\end{lemma}

\begin{proof}
Let $F$ be a $\sigma$-flag of $\T_0$ that is not a model of $\T_1$, and let
$M$ be its underlying unlabelled structure.  Since $\T_1$ is universal, $M$ is
not a model of $\T_1$.  The downward average $\avg{F}{\sigma}$ is a
non-negative scalar multiple of $M$ as an element of $\A^0[\T_0]$.  Since
$\phi_0\in Q_0$, we have $\phi_0(M)=0$, and hence
$\phi_0(\avg{F}{\sigma})=0$.  By \eqref{eq:extension-expectation},
\[
  \E[\phi^\sigma(F)]=0.
\]
But $\phi^\sigma(F)\ge0$ pointwise.  Hence $\phi^\sigma(F)=0$ almost surely.
Intersecting the resulting probability-one events over the countable set of
forbidden $\sigma$-flags shows that $\phi^\sigma\in Q_\sigma$ almost surely.
\end{proof}

\section{The support-closure criterion}

The following definition is the main conceptual point of this paper.

\begin{definition}[Root-planting set and root-plantability]\label{def:root-planting}
Fix a non-degenerate type $\sigma$ of $\T_1$.  Define
\[
  S_\sigma
  :=
  \overline{
  \bigcup_{\substack{\phi_0\in Q_0\\ \phi_0(\brak\sigma)>0}}
  \supp(\Ext_\sigma(\phi_0))}
  \subseteq X_\sigma,
\]
where the closure is taken in the compact ambient space
$X_\sigma=\Homp(\A^\sigma[\T_0],\R)$.
We say that the triple $(\T_0,\T_1,\sigma)$ is \emph{root-plantable} if
\[
  S_\sigma=Q_\sigma.
\]
\end{definition}

By Lemma~\ref{lem:support-passes-general}, one always has $S_\sigma\subseteq Q_\sigma$.
The point of root-plantability is the reverse inclusion: every constrained
labelled homomorphism should be approximable by support points of random
extensions from constrained unlabelled homomorphisms.

\begin{theorem}[Support-closure criterion]\label{thm:support-criterion}
Let $\T_1$ be a universal strengthening of $\T_0$, and let $\sigma$ be a non-degenerate type of $\T_1$.  For $f\in\A^\sigma[\T_0]$, consider the following two
conditions.
\begin{enumerate}[label=\textup{(\alph*)}, leftmargin=2.8em]
  \item\label{cond:quotient}
  Quotient semantics:
  \[
    \chi(f)\ge0\qquad\text{for every }\chi\in Q_\sigma.
  \]
  Equivalently, $\psi(\qmap_\sigma f)\ge0$ for every
  $\psi\in\Homp(\A^\sigma[\T_1],\R)$.
  \item\label{cond:ensemble}
  Ensemble semantics:
  for every $\phi_0\in Q_0$ with $\phi_0(\brak\sigma)>0$,
  \[
    \Prob_{\phi^\sigma\sim\Ext_\sigma(\phi_0)}[\phi^\sigma(f)\ge0]=1.
  \]
\end{enumerate}
Then \ref{cond:quotient} always implies \ref{cond:ensemble}.  Moreover,
\ref{cond:quotient} and \ref{cond:ensemble} are equivalent for every
$f\in\A^\sigma[\T_0]$ if and only if $(\T_0,\T_1,\sigma)$ is root-plantable.
\end{theorem}

\begin{proof}
By Lemma~\ref{lem:support-passes-general}, every support appearing in the definition
of $S_\sigma$ is contained in $Q_\sigma$.  Hence non-negativity on $Q_\sigma$
implies non-negativity on all those supports, and by Lemma~\ref{lem:support-as} this
is precisely the implication from quotient semantics to ensemble semantics.

For the converse direction, note first that ensemble semantics for a fixed $f$
is equivalent, again by Lemma~\ref{lem:support-as}, to non-negativity of $f$ on the
union of the supports in Definition~\ref{def:root-planting}.  Since $f$ is continuous on
$X_\sigma$, this is equivalent to non-negativity on $S_\sigma$.
Thus the comparison is simply
\[
  f\ge0\text{ on }S_\sigma
  \qquad\text{versus}\qquad
  f\ge0\text{ on }Q_\sigma.
\]
If $S_\sigma=Q_\sigma$, the two conditions are equivalent for every $f$.

Conversely suppose $S_\sigma\ne Q_\sigma$.  Choose
$\psi\in Q_\sigma\setminus S_\sigma$.  The compact space $X_\sigma$ is a closed subspace of the product
$[0,1]^{\F^\sigma}$, hence is compact metrizable and normal.  The closed sets
$S_\sigma$ and $\{\psi\}$ are disjoint subsets of the closed space $Q_\sigma$.
By Urysohn's lemma there is a continuous function $h\in C(Q_\sigma)$ such that
$h=1$ on $S_\sigma$ and $h(\psi)=-1$.  By Tietze's extension theorem, extend
$h$ to a continuous function $H\in C(X_\sigma)$.  By the Stone-Weierstrass
theorem in the form used by Razborov \cite[Proposition 3.7]{Razborov2007},
$\A^\sigma[\T_0]$ is uniformly dense in $C(X_\sigma)$.  Choose
$f\in\A^\sigma[\T_0]$ with $\|f-H\|_\infty<1/3$.  Then $f>0$ on
$S_\sigma$, so ensemble semantics holds for $f$, while $f(\psi)<0$, so quotient
semantics fails.  Thus equivalence for all $f$ forces $S_\sigma=Q_\sigma$.
\end{proof}

\section{Clone-closed graph classes are root-plantable}

We now give a checkable sufficient condition for root-plantability for graph
classes, built on an independent blow-up construction.  Let
$\T_0=\T_{\mathrm{Graph}}$, and let $\K$ be a hereditary class of finite simple
graphs.  Equivalently, $\K$ is the class of finite models of the universal theory
$\T_\K$ obtained from graph theory by forbidding all induced graphs not in $\K$.

\begin{definition}[Independent blow-up]
Let $G$ be a finite graph and let $m_v\ge1$ be an integer for each
$v\in V(G)$.  The independent blow-up $G^{\mathbf m}$ is obtained by replacing
each vertex $v$ by an independent clone class $C_v$ of size $m_v$, and by putting
all edges between $C_v$ and $C_w$ whenever $vw\in E(G)$.
\end{definition}

\begin{definition}[Clone-closed hereditary class]
A hereditary graph class $\K$ is \emph{clone-closed} if every independent
blow-up of every graph in $\K$ again belongs to $\K$.
\end{definition}

\begin{example}
For every $r\ge3$, the class of $K_r$-free graphs is clone-closed.  A clique in
an independent blow-up uses at most one vertex from each clone class, and hence
projects to a clique of the same size in the original graph.  In particular, the
triangle-free class is clone-closed.
\end{example}

\subsection{The planted blow-up estimate}

Let $F=(G,\theta)$ be a $\sigma$-flag with $|\sigma|=k$ and $|G|=n>k$.  Fix a
parameter $0<\lambda\le 1/2$.  If $k>0$, assign weights
\[
  w_{\theta(i)}=\frac{\lambda}{k}\quad(1\le i\le k),
  \qquad
  w_v=\frac{1-\lambda}{n-k}\quad(v\notin\operatorname{im}\theta).
\]
If $k=0$, set $w_v=1/n$ for all $v\in V(G)$.  For a large integer $N$, let
$B_N(G,\theta,\lambda)$ be an independent blow-up of $G$ whose clone class
$C_v$ has size approximately $w_vN$, and put
\[
  \operatorname{err}_N
  :=
  \sum_{v\in V(G)}\left|\frac{|C_v|}{N}-w_v\right|.
\]
The error tends to zero along suitable choices of integer cluster sizes.

\begin{lemma}[Planted blow-up approximation]\label{lem:planted-estimate}
Let $m\ge k$ be fixed.  There is a constant $C_m$, depending only on $m$, such
that the following holds for every $0<\lambda\le 1/2$.  Let $F_0$ be any $\sigma$-flag with at most $m$
vertices, let $(G,\theta)$ be a $\sigma$-flag on $n>k$ vertices, and let
$B_N=B_N(G,\theta,\lambda)$.  If
$\widehat\theta:\sigma\to B_N$ satisfies
$\widehat\theta(i)\in C_{\theta(i)}$ for all $i$, then
\begin{equation}\label{eq:planted-estimate}
  \left|p(F_0,(B_N,\widehat\theta))-p(F_0,(G,\theta))\right|
  \le
  C_m\left(\lambda+\frac1{n-k}+\operatorname{err}_N\right).
\end{equation}
Consequently, for every fixed $f\in\A^\sigma[\T_{\mathrm{Graph}}]$ there is a
constant $C_f$ such that
\begin{equation}\label{eq:planted-linear-estimate}
  \left|p(f,(B_N,\widehat\theta))-p(f,(G,\theta))\right|
  \le
  C_f\left(\lambda+\frac1{n-k}+\operatorname{err}_N\right).
\end{equation}
\end{lemma}

\begin{proof}
It suffices to prove \eqref{eq:planted-estimate}; the estimate for $f$ follows
by linearity because $f$ is a finite linear combination of flags.

Let $\ell=|F_0|\le m$.  To compute $p(F_0,(B_N,\widehat\theta))$, choose
$\ell-k$ additional unlabelled vertices uniformly from $B_N$, avoiding the
labelled vertices.  Call the sample good if none of these additional vertices
lies in a planted labelled clone class $C_{\theta(1)},\ldots,C_{\theta(k)}$ and
no two of them lie in the same non-labelled clone class.  The probability of
hitting a planted labelled clone class is at most
$(\ell-k)\lambda+O_m(\operatorname{err}_N)$.  Conditional on avoiding the planted clone classes, the normalised target weights
of the non-labelled clone classes are equal, up to $O_m(\operatorname{err}_N)$,
because the total non-labelled mass is $1-\lambda\ge1/2$.  Therefore the
probability that two sampled vertices
fall in the same non-labelled clone class is at most
$\binom{\ell-k}{2}/(n-k)+O_m(\operatorname{err}_N)$.  Hence the bad probability is
bounded by the right-hand side of \eqref{eq:planted-estimate}, after enlarging
$C_m$.

On the good event, projecting every sampled vertex of $B_N$ to its base vertex
in $G$ gives a uniformly chosen $(\ell-k)$-subset of
$V(G)\setminus\operatorname{im}\theta$, up to total variation error
$O_m(\operatorname{err}_N)$.  The induced labelled graph in the blow-up is then
exactly the induced labelled graph in $(G,\theta)$ on the projected vertices,
because distinct clone classes have precisely the adjacencies of their base
vertices.  Thus the two flag probabilities differ only through the bad event and
rounding error, which proves the claim.
\end{proof}

\begin{lemma}[Positive probability of the planted root]\label{lem:planted-mass}
With the notation above, for all sufficiently large $N$, the finite random
extension of the unlabelled graph $B_N(G,\theta,\lambda)$ to type $\sigma$ gives
probability at least
\[
  c_{k,\lambda}:=
  \begin{cases}
    1, & k=0,\\[1mm]
    \left(\dfrac{\lambda}{2k}\right)^k, & k>0,
  \end{cases}
\]
to the event that the random labelled embedding
$\widehat\theta:\sigma\to B_N(G,\theta,\lambda)$ satisfies
$\widehat\theta(i)\in C_{\theta(i)}$ for all $i=1,\ldots,k$.
\end{lemma}

\begin{proof}
For $k=0$ there is nothing to prove.  Assume $k>0$.  Every ordered choice of
vertices $\widehat\theta(i)\in C_{\theta(i)}$ is a model embedding of $\sigma$,
because $\theta$ is a model embedding into $G$ and the blow-up preserves the
relations between distinct base vertices.  For $N$ sufficiently large,
$|C_{\theta(i)}|\ge(\lambda/(2k))N$ for all $i$.  Thus the number of planted
ordered embeddings is at least $(\lambda/(2k))^kN^k$, while the total number of
ordered embeddings of $\sigma$ into $B_N$ is at most $N^k$.
\end{proof}

\subsection{Root-plantability theorem}

\begin{theorem}[Clone-closed classes are root-plantable]\label{thm:clone-root-plantable}
Let $\K$ be a clone-closed hereditary graph class, let
$\T_1=\T_\K$, and let $\sigma$ be a non-degenerate type of $\T_\K$.  Then
\[
  S_\sigma=Q_\sigma.
\]
Consequently, quotient semantics and ensemble semantics are equivalent for every
$f\in\A^\sigma[\T_{\mathrm{Graph}}]$.
\end{theorem}

\begin{proof}
The inclusion $S_\sigma\subseteq Q_\sigma$ is Lemma~\ref{lem:support-passes-general}.
We prove the reverse inclusion.  Let $\psi\in Q_\sigma$.  We must show that every
neighbourhood of $\psi$ in $X_\sigma$ intersects the support of some
$\Ext_\sigma(\phi_0)$ with $\phi_0\in Q_0$ and $\phi_0(\brak\sigma)>0$.

Let a basic neighbourhood of $\psi$ be specified by flags
$F_1,\ldots,F_s$ and a tolerance $\eps>0$:
\[
  U=
  \{\chi\in X_\sigma:
    |\chi(F_i)-\psi(F_i)|<\eps\text{ for all }i\}.
\]
By the representation theorem in the constrained theory $\T_\K$, choose a
convergent sequence of $\sigma$-flags
$(G_t,\theta_t)$ with $G_t\in\K$, $|G_t|\to\infty$, and
$p(\cdot,(G_t,\theta_t))\to\psi$.  Let $m=\max_i |F_i|$.
Choose $0<\lambda\le 1/2$ so small that $C_m\lambda<\eps/10$, where $C_m$ is the
constant from Lemma~\ref{lem:planted-estimate}.  Then fix one $t$ large enough that
\[
  |p(F_i,(G_t,\theta_t))-\psi(F_i)|<\eps/10
  \quad\text{for all }i,
  \qquad
  \frac{C_m}{|G_t|-|\sigma|}<\eps/10.
\]
Finally take a sequence $N_j\to\infty$ with rounding errors tending to zero and, for all large $j$, satisfying $C_m\operatorname{err}_{N_j}<\eps/10$; put
\[
  B_j:=B_{N_j}(G_t,\theta_t,\lambda).
\]
Since $G_t\in\K$ and $\K$ is clone-closed, every $B_j$ lies in $\K$.  Passing to
a subsequence if necessary, the unlabelled graphs $B_j$ converge to some
$\phi_0\in Q_0$.  The counting argument in Lemma~\ref{lem:planted-mass} gives a positive lower bound, independent of $j$, on the density of embeddings of $\sigma$ in $B_j$; hence $\phi_0(\brak\sigma)>0$.

Let $Y_\sigma:=[0,1]^{\F^\sigma[\T_{\mathrm{Graph}}]}$, so that
$X_\sigma\subseteq Y_\sigma$ is a closed subspace.  Let $P_j$ be the finite
random-root distribution on $Y_\sigma$ obtained by choosing a random embedding
of $\sigma$ into $B_j$ and recording all rooted flag densities.  These measures
converge weakly to $\Ext_\sigma(\phi_0)$, viewed as a probability measure on
$Y_\sigma$ concentrated on $X_\sigma$.

Consider the closed cylinder in the ambient product space
\[
  \widetilde C=
  \{x\in Y_\sigma:
    |x(F_i)-\psi(F_i)|\le\eps/2\text{ for all }i\}.
\]
Put $C:=\widetilde C\cap X_\sigma$.  Then $C\subseteq U$.  For all sufficiently
large $j$, Lemma~\ref{lem:planted-estimate} implies that every planted root lies
in $\widetilde C$.  By Lemma~\ref{lem:planted-mass},
\[
  P_j(\widetilde C)\ge c_{|\sigma|,\lambda}>0
\]
for all sufficiently large $j$.  Since $\widetilde C$ is closed in $Y_\sigma$, the
closed-set form of the Portmanteau theorem gives
\[
  \Ext_\sigma(\phi_0)(\widetilde C)
  \ge
  \limsup_{j\to\infty} P_j(\widetilde C)
  \ge c_{|\sigma|,\lambda}>0.
\]
Since $\Ext_\sigma(\phi_0)$ is concentrated on $X_\sigma$, this says
$\Ext_\sigma(\phi_0)(C)>0$.  Hence the support of $\Ext_\sigma(\phi_0)$ intersects
$C\subseteq U$.  Since $U$ was arbitrary, $\psi\in S_\sigma$.
\end{proof}

\begin{corollary}[Clique-free and triangle-free theories]\label{cor:clique-free}
For every $r\ge3$, every non-degenerate $K_r$-free type $\sigma$, and every
$f\in\A^\sigma[\T_{\mathrm{Graph}}]$, quotient semantics and ensemble semantics
are equivalent for the $K_r$-free theory.  In particular, the equivalence holds
for the triangle-free theory.
\end{corollary}

\section{Degenerate classes are not root-plantable}

The clone-closed hypothesis is not merely technical: the correspondence fails for
hereditary classes, and it fails for an entire family, always for the same
structural reason.  The obstruction is \emph{degeneracy}.  If a class is so sparse
that its dense limits carry no edges, then a typical root sees an empty
neighbourhood; yet the quotient theory may still retain an \emph{exceptional}
high-degree root, such as the centre of a star.  The rooted edge flag separates the
two.

Throughout this section $\sigma=1$ is the one-vertex type in graph theory,
$e\in\A^1[\T_{\mathrm{Graph}}]$ is the rooted edge flag (the two-vertex flag in which
the labelled vertex is adjacent to the one unlabelled vertex), and $\rho\in\A^0$ is
the unlabelled edge flag; recall $\avg{e}{1}=\rho$ and $\brak1=1$.

\begin{definition}[Edge-degenerate class]\label{def:edge-degenerate}
A hereditary graph class $\K$ is \emph{edge-degenerate} if every constrained
unlabelled limit has edge density zero, that is, $\phi_0(\rho)=0$ for all
$\phi_0\in Q_0$.  Equivalently, the maximum number of edges of an $n$-vertex graph
in $\K$ is $o(n^2)$.
\end{definition}

\begin{theorem}[Degeneracy obstruction]\label{thm:degenerate-obstruction}
Let $\K$ be a hereditary graph class that is edge-degenerate and contains
arbitrarily large stars.  Then $(\T_{\mathrm{Graph}},\T_\K,\sigma{=}1)$ is not
root-plantable.  Concretely, $f:=-e\in\A^1[\T_{\mathrm{Graph}}]$ satisfies ensemble
semantics but not quotient semantics, and in fact
\[
  S_1\ \subseteq\ \{\chi\in Q_1:\chi(e)=0\}\ \subsetneq\ Q_1 .
\]
\end{theorem}

\begin{proof}
\emph{Ensemble semantics holds for $-e$.}  Let $\phi_0\in Q_0$.  Since $\brak1=1$ we
have $\phi_0(\brak1)=1>0$, so $\Ext_1(\phi_0)$ is defined.  As $\avg{e}{1}=\rho$,
\eqref{eq:extension-expectation} gives
\[
  \E_{\phi^1\sim\Ext_1(\phi_0)}[\phi^1(e)]
  =\frac{\phi_0(\rho)}{\phi_0(\brak1)}=\phi_0(\rho)=0,
\]
using edge-degeneracy.  Because $\phi^1(e)\ge0$ pointwise, $\phi^1(e)=0$ almost
surely, hence $\phi^1(-e)=0\ge0$ almost surely; ensemble semantics holds.  Moreover,
by Lemma~\ref{lem:support-as} every support point $\chi$ of every $\Ext_1(\phi_0)$
satisfies $\chi(e)=0$, and since $\chi\mapsto\chi(e)$ is continuous this gives the
inclusion $S_1\subseteq\{\chi\in Q_1:\chi(e)=0\}$.

\emph{Quotient semantics fails for $-e$.}  Let $S_N=K_{1,N-1}$ be the star on $N$
vertices, rooted at its centre; by hypothesis $S_N\in\K$ for arbitrarily large $N$.
Every unlabelled vertex is adjacent to the centre, so $p(e,S_N)=1$ for all $N$.  By
compactness a subsequence converges to some $\psi\in Q_1$ with $\psi(e)=1$, whence
$\psi(-e)=-1<0$.  Thus quotient semantics fails, and $\psi\in Q_1$ witnesses
$\{\chi\in Q_1:\chi(e)=0\}\subsetneq Q_1$; in particular $\psi\in Q_1\setminus S_1$,
so $S_1\ne Q_1$.
\end{proof}

\begin{remark}
The hypothesis ``contains arbitrarily large stars'' is only used to exhibit a
quotient root of positive rooted edge density; it may be replaced by the existence
of any sequence of rooted graphs in $\K$ whose rooted edge density stays bounded
away from $0$.
\end{remark}

\begin{remark}[The obstruction]
The centre of a star is an exceptional labelled vertex.  Quotient semantics can see
it, because labelled limits may root at exceptional vertices.  Ensemble semantics
instead roots a constrained unlabelled limit at a \emph{typical} vertex; in a
degenerate class the edge density is zero, so a typical vertex has rooted edge
density zero.  The high-degree root cannot be planted with positive probability
without leaving the class: duplicating the centre and one neighbour already produces
a $K_{2,2}$, and more generally a denser local structure than a degenerate class can
sustain.
\end{remark}

\subsection{Examples}

The hypotheses of Theorem~\ref{thm:degenerate-obstruction} are met by many familiar
sparse classes.  We first record the original $C_4$-free case, where edge-degeneracy
is the classical K\H{o}v\'ari--S\'os--Tur\'an bound.

Let
\[
  \K_{C_4}:=\{G:\text{$G$ contains no copy of $C_4$ as a not-necessarily-induced subgraph}\}.
\]
This is hereditary; in the induced-subgraph language of flag algebras it is
axiomatised by forbidding the three four-vertex graphs that contain a $C_4$ as a
subgraph: $C_4$, $K_4$ minus an edge, and $K_4$.  It is not clone-closed: the
independent blow-up of $K_2$ into two clone classes of size $2$ is $K_{2,2}=C_4$.

\begin{lemma}[Dense \texorpdfstring{$C_4$}{C4}-free limits have edge density zero]\label{lem:c4-edge-zero}
$\K_{C_4}$ is edge-degenerate: if $\phi_0\in Q_0$ for the $C_4$-free theory, then
$\phi_0(\rho)=0$.
\end{lemma}

\begin{proof}
By the representation theorem in the constrained theory, $\phi_0$ is the limit
of a convergent sequence of finite $C_4$-free graphs $G_n$ with $|G_n|\to\infty$.
We show that their edge densities tend to zero.

Let $G$ be a $C_4$-free graph on $N$ vertices, with degrees $d(v)$.  Any two
vertices have at most one common neighbour; otherwise two common neighbours and
the two original vertices form a $C_4$.  Therefore
\[
  \sum_{v\in V(G)} \binom{d(v)}{2}
  \le
  \binom{N}{2}.
\]
By convexity,
\[
  N\binom{\overline d}{2}
  \le
  \binom{N}{2},
  \qquad
  \overline d:=\frac1N\sum_v d(v).
\]
Thus $\overline d(\overline d-1)\le N-1$, so
$\overline d\le 1+\sqrt N$.  The edge density is
$\overline d/(N-1)$, which tends to zero.  Hence
$\phi_0(\rho)=0$.
\end{proof}

\begin{corollary}[\texorpdfstring{$C_4$}{C4}-free counterexample]\label{cor:c4-counterexample}
$\K_{C_4}$ is edge-degenerate (Lemma~\ref{lem:c4-edge-zero}) and contains all stars,
so by Theorem~\ref{thm:degenerate-obstruction} it is not root-plantable at
$\sigma=1$: the element $-e$ satisfies ensemble semantics but not quotient
semantics.
\end{corollary}

\begin{corollary}[A family of degenerate obstructions]\label{cor:degenerate-family}
Each of the following hereditary classes is edge-degenerate and contains all stars,
hence is not root-plantable at the one-vertex type:
\begin{enumerate}[label=\textup{(\roman*)}, leftmargin=2.5em]
  \item $K_{s,t}$-subgraph-free graphs for $2\le s\le t$, edge-degenerate by the
        K\H{o}v\'ari--S\'os--Tur\'an bound $\mathrm{ex}(n,K_{s,t})=O(n^{2-1/s})$
        \cite{KovariSosTuran} (a star contains no $K_{2,2}$);
  \item $C_{2k}$-subgraph-free graphs for $k\ge2$, edge-degenerate by the
        Bondy--Simonovits bound $\mathrm{ex}(n,C_{2k})=O(n^{1+1/k})$
        \cite{BondySimonovits};
  \item forests, and more generally any class of bounded average degree containing
        all stars ($O(n)$ edges);
  \item planar graphs ($\le 3n-6$ edges).
\end{enumerate}
\end{corollary}

\subsection{Complementation: a dense obstruction}

One might hope that the failures above are an artefact of \emph{sparsity}, and that
density rescues root-plantability.  It does not.  Root-plantability is preserved by
complementation, which immediately turns each sparse obstruction into a dense one.

\begin{lemma}[Complementation invariance]\label{lem:complementation}
For a hereditary graph class $\K$ let $\overline\K=\{\overline G:G\in\K\}$, which is
again hereditary, and for a type $\sigma$ let $\overline\sigma$ be its complement.
Then $(\T_{\mathrm{Graph}},\T_\K,\sigma)$ is root-plantable if and only if
$(\T_{\mathrm{Graph}},\T_{\overline\K},\overline\sigma)$ is.  For the one-vertex type
$\sigma=1=\overline\sigma$, the class $\K$ is root-plantable at $1$ if and only if
$\overline\K$ is.
\end{lemma}

\begin{proof}
Vertex-wise complementation $G\mapsto\overline G$ is an involution on finite graphs
that takes induced-subgraph densities to those of the complements, hence extends to
an isomorphism of flag algebras
$\A^\sigma[\T_{\mathrm{Graph}}]\cong\A^{\overline\sigma}[\T_{\mathrm{Graph}}]$ carrying
$\K$-flags to $\overline\K$-flags.  It therefore induces a homeomorphism
$X_\sigma\cong X_{\overline\sigma}$ taking $Q_\sigma$ for $\K$ onto $Q_{\overline\sigma}$
for $\overline\K$.  Complementation also commutes with the random root: a uniformly
random embedding of $\sigma$ into $G$ is the same embedding of $\overline\sigma$ into
$\overline G$, so it carries each $\Ext_\sigma(\phi_0)$ to the corresponding
extension for $\overline\K$, and hence $S_\sigma$ for $\K$ onto $S_{\overline\sigma}$
for $\overline\K$.  Thus $S_\sigma=Q_\sigma$ holds for $\K$ if and only if it holds
for $\overline\K$.
\end{proof}

Call $\K$ \emph{co-edge-degenerate} if $\overline\K$ is edge-degenerate; equivalently
every $\phi_0\in Q_0$ has $\phi_0(\rho)=1$, equivalently every $n$-vertex graph in
$\K$ has at least $\binom n2-o(n^2)$ edges.  Such classes are as dense as possible,
with edge density tending to $1$, yet they fail just as the sparse ones do.

\begin{corollary}[A dense obstruction]\label{cor:codegenerate}
Let $\K$ be hereditary, co-edge-degenerate, and contain all co-stars
$\overline{K_{1,N-1}}=K_{N-1}\sqcup K_1$.  Then $\K$ is not root-plantable at
$\sigma=1$, witnessed by $f:=e-1_1$ (with $1_1$ the unit of $\A^1$).  In particular the complements of the classes in
Corollary~\ref{cor:degenerate-family}---complement-of-$C_4$-free,
complement-of-$K_{s,t}$-free, complement-of-even-cycle-free, and
complement-of-planar graphs---are dense hereditary classes that are not
root-plantable.
\end{corollary}

\begin{proof}
By Lemma~\ref{lem:complementation} it suffices that $\overline\K$ be not
root-plantable at $1$, and this is Theorem~\ref{thm:degenerate-obstruction}:
$\overline\K$ is edge-degenerate and contains all stars $K_{1,N-1}$.  The witness
$-e$ for $\overline\K$ pulls back to $-(1_1-e)=e-1_1$ for $\K$.  Directly:
$\avg{e}{1}=\rho$ gives $\E[\phi^1(e)]=\phi_0(\rho)=1$, and $\phi^1(e)\le1$ forces
$\phi^1(e)=1$ almost surely, so $e-1_1$ satisfies ensemble semantics; while the
co-star rooted at its isolated vertex gives a $\psi\in Q_1$ with $\psi(e)=0$, whence
$\psi(e-1_1)=-1<0$ and quotient semantics fails.
\end{proof}

\begin{remark}[The dichotomy: density is not the dividing line]
Theorem~\ref{thm:degenerate-obstruction} and Corollary~\ref{cor:codegenerate}
together show that root-plantability is governed neither by heredity nor by density:
edge-degenerate classes (rooted edge density pinned to $0$) and co-edge-degenerate
classes (pinned to $1$) both fail, for the same reason.  What separates a typical
root from the exceptional one is that the rooted edge density is \emph{pinned} to a
boundary value of $[0,1]$ across all constrained unlabelled limits, while the
quotient theory still contains a root attaining the opposite boundary.  The
governing invariant is this boundary pinning, not sparsity or density.

This refines, rather than closes, the problem of characterising root-plantability.
The $C_5$-free graphs remain the natural open test case.  Their edge density is
\emph{not} pinned---the class contains graphs of every density in $[0,\tfrac12]$, so
no boundary-pinning obstruction applies---yet the class is not clone-closed (the
blow-up of a triangle into the octahedron $K_{2,2,2}$ contains a $C_5$).  Whether it
is root-plantable thus tests whether the absence of boundary pinning is by itself
sufficient, the frontier left open between Theorem~\ref{thm:clone-root-plantable}
and the two obstructions.
\end{remark}

\subsection{The general obstruction and a conjecture}

Both obstructions are instances of one mechanism: a rooted density pinned almost
surely to a single value across all constrained limits, while the quotient escapes
it.

\begin{theorem}[Pinning obstruction]\label{thm:pinning}
Let $\sigma$ be a non-degenerate type, let $g$ be a $\sigma$-flag, and let
$c\in[0,1]$.  Suppose the rooted density of $g$ is almost surely pinned to $c$ across
all constrained unlabelled limits: for every $\phi_0\in Q_0$ with
$\phi_0(\brak\sigma)>0$, one has $\phi^\sigma(g)=c$ almost surely under
$\Ext_\sigma(\phi_0)$.  If the quotient theory contains a root attaining a different
value---some $\psi\in Q_\sigma$ with $\psi(g)\ne c$---then $(\T_0,\T_1,\sigma)$ is not
root-plantable.
\end{theorem}

\begin{proof}
Both $g-c\,1_\sigma$ and $c\,1_\sigma-g$ have $\phi^\sigma$-value $0$ almost surely
under every admissible $\Ext_\sigma(\phi_0)$, so each is ensemble-nonnegative and,
by Lemma~\ref{lem:support-as} and continuity, nonnegative on $S_\sigma$.  Since
$\psi(g)\ne c$, one of the two is strictly negative at $\psi$; that
$\psi\in Q_\sigma$ therefore lies outside $S_\sigma$, so $S_\sigma\ne Q_\sigma$.
\end{proof}

\begin{remark}[Endpoints are automatic and checkable]
At the endpoints $c\in\{0,1\}$ the almost-sure pinning follows from a mere
expectation condition, since a $[0,1]$-valued random variable with mean $0$
(respectively $1$) is almost surely $0$ (respectively $1$).  Concretely, by
\eqref{eq:extension-expectation}, $\phi_0(\avg{g}{\sigma})=0$ for all admissible
$\phi_0$ gives the case $c=0$ (witness $-g$), and
$\phi_0(\avg{g}{\sigma})=\phi_0(\brak\sigma)$ gives $c=1$ (witness $g-1_\sigma$).  The
degeneracy obstruction (Theorem~\ref{thm:degenerate-obstruction}) and its dense
counterpart (Corollary~\ref{cor:codegenerate}) are exactly these, with $\sigma=1$ and
$g=e$.  An interior value $c\in(0,1)$ instead requires the almost-sure pinning as a
genuine hypothesis.
\end{remark}

\begin{conjecture}[Characterisation; tentative]\label{conj:characterisation}
$(\T_{\mathrm{Graph}},\T_\K,\sigma)$ is root-plantable if and only if it admits no
pinning obstruction: there is no $\sigma$-flag $g$ and constant $c$ with
$\phi^\sigma(g)=c$ almost surely for all constrained unlabelled limits while some
$\psi\in Q_\sigma$ has $\psi(g)\ne c$.
\end{conjecture}

The forward implication is Theorem~\ref{thm:pinning}; the conjecture asserts that the
absence of such single-flag pinning is already sufficient.  It is tentative:
root-plantability could conceivably fail through a \emph{joint} constraint on several
flags with no single flag pinned, in which case a richer notion of obstruction would
be needed.

\section{Toward the \texorpdfstring{$C_5$}{C5}-free case}

Conjecture~\ref{conj:characterisation} predicts that the $C_5$-free graphs---dense,
not clone-closed, and, as we now confirm, free of any pinning obstruction from the
edge or triangle---are root-plantable.  We record the structural evidence;
throughout, $C_5$-free again means containing no $C_5$ as a not-necessarily-induced
subgraph (the densest such graphs being complete bipartite).

The key fact is a strong local constraint.

\begin{lemma}[Neighbourhoods in $C_5$-free graphs]\label{lem:c5-nbhd}
If $G$ is $C_5$-free and $v\in V(G)$, then $G[N(v)]$ contains no path on four
vertices as a subgraph; equivalently, each neighbourhood induces a disjoint union of
stars and triangles.
\end{lemma}

\begin{proof}
A path $a$-$x$-$y$-$b$ inside $N(v)$ together with $v$ forms the $5$-cycle
$v$-$a$-$x$-$y$-$b$-$v$.  Thus $G[N(v)]$ is $P_4$-subgraph-free, and the connected
$P_4$-subgraph-free graphs are exactly the stars $K_{1,t}$ and the triangle.
\end{proof}

\begin{corollary}[No pinning obstruction for $C_5$-free]\label{cor:c5-no-pin}
The $C_5$-free theory has no pinning obstruction at $\sigma=1$ from the edge or the
rooted-triangle flag.
\end{corollary}

\begin{proof}
A disjoint union of stars and triangles on at most $\deg(v)$ vertices has at most
$\deg(v)$ edges, so the number of triangles through $v$ is at most $\deg(v)\le n$,
and the rooted-triangle density at every vertex is $O(1/n)$.  Hence the
rooted-triangle flag has density $0$ on \emph{all} of $Q_1$: it is pinned at $0$, but
the quotient attains only $0$, so no obstruction arises.  The edge flag is not
pinned at all, as $C_5$-free graphs realise every edge density in $[0,\tfrac12]$
(the empty graph and the complete bipartite graphs).
\end{proof}

Finally, blow-ups detect exactly the triangles.

\begin{lemma}\label{lem:c5-blowup}
Every independent blow-up of a $C_5$-free graph $G$ is again $C_5$-free if and only
if $G$ is triangle-free.
\end{lemma}

\begin{proof}
A $C_5$ in a blow-up projects to a closed walk of length $5$ in $G$, and conversely
any closed walk of length $5$ in $G$ lifts to a $C_5$ in a blow-up with classes of
size $\ge2$.  Since $G$ is $C_5$-free, such a walk---being of odd length $5$, hence
containing an odd cycle of length $3$ or $5$---exists if and only if $G$ has a
triangle.
\end{proof}

Consequently the triangle-free $C_5$-free graphs form a clone-closed class, so by
Theorem~\ref{thm:clone-root-plantable} every $C_5$-free rooted limit realised by
triangle-free graphs lies in $S_1$; and by Lemma~\ref{lem:c5-nbhd} the
neighbourhood of a positive-degree root is asymptotically independent (its induced
edges have vanishing density), so dense $C_5$-free roots are locally bipartite.
These facts leave no candidate obstruction, and we conjecture:

\begin{conjecture}\label{conj:c5}
The $C_5$-free theory is root-plantable at every non-degenerate type; in particular
Conjecture~\ref{conj:characterisation} holds for it.
\end{conjecture}

\section{What this says about strengthening by another theory \texorpdfstring{$\T'$}{T'}}

The discussion above gives a clean answer to the proposed generalisation from a
hereditary class $\K$ to a conjunction $\T\wedge\T'$.

\subsection{When \texorpdfstring{$\T'$}{T'} is universal}

If $\T'$ is universal and uses the same relational language as $\T$, then
$\T_1:=\T\wedge\T'$ is again a universal theory.  The quotient flag algebra
$\A^\sigma[\T_1]$ is well-defined, and the support-closure criterion
Theorem~\ref{thm:support-criterion} applies verbatim.  Thus the next theorem in this
direction should not claim equivalence for all universal strengthenings.  The
degeneracy obstruction (Theorem~\ref{thm:degenerate-obstruction}), and already the
$C_4$-free instance, shows that this is false.  The correct statement is:
\[
  \text{quotient semantics equals ensemble semantics}
  \quad\Longleftrightarrow\quad
  \text{root-plantability.}
\]
The remaining mathematical work is therefore to find useful, checkable sufficient
conditions for root-plantability.

Clone-closure is one such sufficient condition for graph classes.  Other possible
conditions worth investigating are closure under suitable substitutions,
true-twin or false-twin duplications, or more model-theoretic amalgamation and
duplication properties.

\subsection{When \texorpdfstring{$\T'$}{T'} is not universal}

If $\T'$ is not universal, then $\T\wedge\T'$ is generally not closed under
induced submodels.  This breaks the basic starting point of Razborov's flag
algebra construction, which assumes that every induced substructure of a finite
model is again a model of the theory.  In this case there may be no canonical
quotient algebra $\A^\sigma[\T\wedge\T']$ of the usual flag-algebra kind.

One can still study constraints imposed on ambient homomorphisms, for example
conditions of the form
\[
  \phi(g_i)=0,
  \qquad
  \phi(h_j)\ge0.
\]
But this is better viewed as \emph{relative semantics inside the ambient flag
algebra}, not as quotient semantics of a new universal theory.  It may be an
interesting direction, but it should be separated from the quotient/ensemble
problem considered here.

\section{Recommended next research targets}

The results suggest three natural next steps.

\begin{enumerate}[label=\textup{(\arabic*)}, leftmargin=2.8em]
  \item \textbf{Develop root-plantability as the main theorem.}
  The support-closure criterion is exact and abstract.  It should become the
  organising statement.  Concrete graph-theoretic hypotheses should then be
  presented as sufficient conditions for root-plantability.

  \item \textbf{Classify duplication operations that imply root-plantability.}
  Clone-closed hereditary classes are only the first example.  The proof uses
  independent duplication of labelled vertices into positive-mass clone classes.
  A broader theorem may allow more general local replacements or substitutions,
  as long as the labelled finite model can be amplified into positive measure
  without leaving the constrained class.

  \item \textbf{Use degeneracy as the canonical obstruction.}
  The centre-rooted star shows exactly what goes wrong: a labelled quotient
  homomorphism may live at an exceptional vertex that cannot be made typical.
  Theorem~\ref{thm:degenerate-obstruction} turns this into a general phenomenon---%
  every edge-degenerate hereditary class with stars fails, including
  $K_{s,t}$-free, even-cycle-free, planar, and bounded-degeneracy classes---and by
  complementation (Corollary~\ref{cor:codegenerate}) the same holds for the
  \emph{dense} co-degenerate classes, so density is not the dividing line.  Both are
  instances of the single pinning obstruction (Theorem~\ref{thm:pinning}).  The
  remaining problem is to settle Conjecture~\ref{conj:characterisation}---that the
  absence of any pinning obstruction is sufficient---for which $C_5$-free graphs
  (Conjecture~\ref{conj:c5}) are the natural first test; Lemmas~\ref{lem:c5-nbhd}
  and~\ref{lem:c5-blowup} reduce that case to planting dense, locally bipartite
  roots.
\end{enumerate}

\begin{thebibliography}{9}

\bibitem{Razborov2007}
A.~A.~Razborov.
\newblock Flag algebras.
\newblock \emph{Journal of Symbolic Logic}, 72(4):1239--1282, 2007.

\bibitem{KovariSosTuran}
T.~K\H{o}v\'ari, V.~T.~S\'os, and P.~Tur\'an.
\newblock On a problem of K. Zarankiewicz.
\newblock \emph{Colloquium Mathematicae}, 3:50--57, 1954.

\bibitem{BondySimonovits}
J.~A.~Bondy and M.~Simonovits.
\newblock Cycles of even length in graphs.
\newblock \emph{Journal of Combinatorial Theory, Series B}, 16(2):97--105, 1974.

\end{thebibliography}

\end{document}
